\documentclass[10pt,twocolumn,letterpaper]{article}

% Replace this preamble with the official CVPR template preamble when exporting.
\usepackage[margin=0.75in]{geometry}
\usepackage{times}
\usepackage{graphicx}
\usepackage{amsmath}
\usepackage{amssymb}
\usepackage{booktabs}
\usepackage{hyperref}

\title{Reason-Then-Retrieve for CoVR-R Composed Video Retrieval}

\author{
MIC Lab\\
}

\begin{document}
\maketitle

\begin{abstract}
We describe a zero-shot system for CoVR-R reason-aware composed video retrieval.
Given a reference video and an edit instruction, the goal is to retrieve the target video from a gallery.
Our method uses a vision-language model to transform the composed query into a target-video representation before retrieval.
The model generates retrieval-oriented gallery descriptions, structured edit reasoning, and edited target-video descriptions.
Dense retrieval is performed by pooling generated-token hidden states into embeddings and ranking candidates by cosine similarity.
We further add a TF-IDF branch over the generated descriptions and combine dense and sparse scores.
The resulting pipeline emphasizes object identity, action order, final state, scene context, and transition cues.
\end{abstract}

\section{Introduction}

Composed video retrieval differs from standard text-video retrieval because the target is specified indirectly.
The query contains a reference video $V_r$ and an edit instruction $E$, while the desired target video $V_t$ must be inferred by applying the edit to the reference.
This setup is challenging when edits modify fine-grained states, action order, hand-object interactions, background context, or camera transitions.

Our system treats the task as a two-step reasoning and retrieval problem.
First, a vision-language model interprets the reference video and the edit instruction to produce a target-oriented description.
Second, both the target-oriented query and the gallery videos are mapped into retrieval features.
We combine hidden-state dense embeddings with sparse lexical matching to benefit from both semantic abstraction and exact visual terms.

\section{Method}

\subsection{Overview}

The pipeline has four stages:

\begin{enumerate}
    \item gallery description and embedding,
    \item query-side edit reasoning and target description,
    \item dense and sparse retrieval,
    \item score fusion.
\end{enumerate}

The system is zero-shot and does not fine-tune the underlying model.
All task adaptation is handled through prompt design, hidden-state pooling, and retrieval-score fusion.

\subsection{Gallery Representation}

For every gallery video, the model generates a compact retrieval-oriented description.
The prompt asks for fields such as the main object, secondary objects, object attributes, surface and scene, action chain, hand interaction, final state, scene transition, camera transition, and distractors.
This structured form helps preserve fine-grained temporal and visual evidence that is often lost in a generic caption.

We extract a dense representation from the generated tokens.
Let $h_i$ denote the hidden state associated with generated token $i$ and $w_i$ its token weight.
The pooled vector is
\begin{equation}
    g = \frac{\sum_i w_i h_i}{\sum_i w_i}.
\end{equation}
Structural field labels and common stop words receive small weights, while content tokens describing objects, actions, and states dominate the pooled representation.
The final gallery vector is L2-normalized.

\subsection{Query Representation}

For each query, the model first observes the reference video and the edit instruction.
It produces structured edit reasoning that records which elements should be preserved, removed, or replaced, and what the target action chain and final state should be.
The model then generates a target-video description conditioned on the reference video, edit instruction, and reasoning trace.

The query embedding is pooled from the generated target-description tokens using the same weighted hidden-state pooling strategy.
Edit-relevant terms extracted from the instruction and reasoning trace receive additional weight.
This biases the representation toward the intended target video rather than the unmodified source.

\subsection{Dense Retrieval}

Dense retrieval ranks candidates using cosine similarity between the query embedding $q$ and each gallery embedding $g$:
\begin{equation}
    s_{\mathrm{dense}}(q,g) = \frac{q^\top g}{\|q\|\|g\|}.
\end{equation}
The ranked list from this branch captures broad semantic compatibility between the inferred target and gallery videos.

\subsection{TF-IDF Retrieval}

Dense retrieval can underweight exact object names, surfaces, and state words.
We therefore add a text branch using TF-IDF over the generated query and gallery descriptions.
The sparse branch uses unigram and bigram features with sublinear term frequency.
It provides a complementary score $s_{\mathrm{tfidf}}$ that often helps when exact visual terms are discriminative.

\subsection{Score Fusion}

Dense and sparse scores are normalized per query and combined:
\begin{equation}
    s_{\mathrm{final}} =
    \alpha \hat{s}_{\mathrm{dense}} +
    (1-\alpha)\hat{s}_{\mathrm{tfidf}}.
\end{equation}
We use split-specific fusion weights selected on development experiments.
The final ranking is obtained by sorting candidates by $s_{\mathrm{final}}$.

\section{Implementation Details}

The main implementation files are:

\begin{itemize}
    \item \texttt{generate\_gallery\_embeddings.py} for gallery description and embedding extraction,
    \item \texttt{run\_retrieval.py} for query reasoning, query embeddings, and dense retrieval,
    \item \texttt{evaluate\_text\_similarity.py} for TF-IDF retrieval,
    \item \texttt{apply\_score\_fusion.py} for dense/sparse score fusion,
    \item \texttt{prompts.py} for the prompt families.
\end{itemize}

Video inputs are processed at a low sampling rate suitable for retrieval.
Generated hidden states are pooled into dense vectors and normalized before cosine search.
The code supports separate WebVid and Something-Something V2 processing, allowing split-specific prompts and fusion settings.

\section{Ablation And Analysis}

We found that dense retrieval and sparse generated-text retrieval are complementary.
Dense embeddings improve semantic matching and tolerate paraphrases, while TF-IDF preserves exact object, surface, and action terms.
Fusing the two branches consistently improved the robustness of the ranking in development experiments.

Prompt design was especially important for action-centric videos.
Descriptions that only summarize the overall scene often miss the temporal order or the final settled state.
Structured prompts that explicitly preserve action chains, end states, scene transitions, and hand-object interactions produced more useful retrieval features.

\section{Observations}

The main error source is often not the retrieval mechanism itself, but the fidelity of the generated text representation.
If the gallery description misses a key action step or the query-side target description keeps removed source content, both dense and sparse retrieval can be misled.
Conversely, when the generated descriptions preserve the edit-critical evidence, simple cosine retrieval plus sparse fusion can be effective.

\section{Limitations}

The approach depends on the quality and consistency of the vision-language model's generated descriptions.
It is also computationally expensive because gallery and query videos require model generation before retrieval.
The score-fusion weights are selected empirically and may need adjustment under a different data distribution.

\section{Conclusion}

We presented a zero-shot reason-then-retrieve system for composed video retrieval.
The method converts a reference-video edit query into a target-video description, extracts hidden-state embeddings from generated tokens, and combines dense retrieval with sparse TF-IDF matching.
The key practical insight is to preserve edit-critical evidence such as object identity, action order, final state, and scene transitions throughout both query and gallery representations.

{\small
\bibliographystyle{ieee_fullname}
\bibliography{references}
}

\end{document}
